\documentclass[aspectratio=169]{beamer}

\graphicspath{{figs}}

\usepackage[T1,T2A]{fontenc}
\usepackage[main=russian,english]{babel}

\usepackage{svg}
\usepackage{cmap}
\usepackage{tikz}
\usetikzlibrary{positioning,arrows.meta}
\usepackage{array}
\usepackage{cancel}
\usepackage{multicol}
\usepackage{booktabs}
\usepackage{csquotes}
\usepackage{amssymb,amsfonts,amsmath}

\usecolortheme{beaver}
\usefonttheme[onlymath]{serif}
\setbeamertemplate{navigation symbols}{}
\setbeamertemplate{footline}[page number]
\setbeamertemplate{blocks}[rounded=true,shadow=true]
\setbeamersize{text margin left=0.5cm,text margin right=0.5cm}
\setbeamerfont{author}{size=\normalsize}
\setbeamerfont{institute}{size=\normalsize}
\setbeamerfont{date}{size=\normalsize}
\setbeamertemplate{enumerate item}{\insertenumlabel)}
\setbeamertemplate{enumerate subitem}{\insertenumlabel)}
\setbeamertemplate{enumerate subsubitem}{\insertenumlabel)}

\DeclareMathOperator*{\argmin}{arg\,min}
\DeclareMathOperator*{\argmax}{arg\,max}

\title{Пространственно-временные характеристики в задаче декодирования мультимодальных многомерных временных рядов}
\author{
    Дорин Даниил Дмитриевич
}
\institute{
    Диссертация на соискание ученой степени\\
    кандидата физико-математических наук\\
    1.2.1 --- Искусственный интеллект и машинное обучение\\
    Научный руководитель: к.ф.-м.н. Грабовой~А.\,В.
}
\date{Москва --- \the\year}

\begin{document}

\begin{frame}
    \thispagestyle{empty}
    \maketitle
\end{frame}

%----------------------------------------------------------------------------------------------------------
\begin{frame}{Задача декодирования мультимодальных временных рядов}
    %% TODO: мотивация, три постановки (лаг, классификация, реконструкция), общая схема
\end{frame}

%----------------------------------------------------------------------------------------------------------
\begin{frame}{Оценка гемодинамического лага}
    %% TODO: линейная авторегрессионная модель, минимум MSE на зрительной коре
\end{frame}

%----------------------------------------------------------------------------------------------------------
\begin{frame}{Пространственно-временная классификация фМРТ}
    %% TODO: маски активности, риманова геометрия, сравнение с нейросетями
\end{frame}

%----------------------------------------------------------------------------------------------------------
\begin{frame}{Мультимодальная реконструкция фМРТ--ЭЭГ}
    %% TODO: контрастивное выравнивание с CLIP, двухстадийная генерация, CLIP-Score
\end{frame}

%----------------------------------------------------------------------------------------------------------
\begin{frame}{Положения, выносимые на защиту}
    %% TODO: из common/characteristic.tex
\end{frame}

\end{document}
